\documentclass[11pt,a4paper]{article}
\usepackage[utf8]{vietnam}
\usepackage{amsmath,amssymb}
\usepackage{tabularx}
\usepackage{tikz}
\usepackage{xcolor}
\usepackage{fancyhdr}
\usepackage[top=1.5cm, bottom=1.5cm, left=1.6cm, right=1.6cm]{geometry}

\setlength{\parskip}{1.5pt}
\setlength{\parindent}{0pt}
\setlength{\abovedisplayskip}{3pt}
\setlength{\belowdisplayskip}{3pt}

\pagestyle{fancy}
\fancyhf{}
\renewcommand{\headrulewidth}{0.4pt}
\lhead{\footnotesize \textbf{CHÙA HƯƠNG SƠN -- ÔN THI VÀO 10}}
\rhead{\footnotesize \textbf{ĐÁP ÁN \& LỜI GIẢI CHI TIẾT (ĐỀ SỐ 10)}}
\cfoot{\thepage}

\begin{document}

\begin{center}
\textbf{\large HƯỚNG DẪN CHẤM VÀ ĐÁP ÁN ĐỀ THI THỬ SỐ 10}\\
\textbf{MÔN: TOÁN 9 -- HỌC KỲ I -- NĂM HỌC 2026 -- 2027}\\
\textit{(Hướng dẫn chấm gồm 02 trang)}
\end{center}

\vspace{1mm}
\noindent
\rule{\linewidth}{0.8pt}

\vspace{1.5mm}
\noindent
\textbf{\large A. ĐÁP ÁN TRẮC NGHIỆM (7,0 ĐIỂM)}

\vspace{1.5mm}
\noindent
\textbf{1. Phần 1: Trắc nghiệm nhiều lựa chọn (3,0 điểm -- Mỗi câu đúng 0,25 điểm)}
\begin{center}
\begin{tabular}{|c|c|c|c|c|c|c|c|c|c|c|c|}
\hline
\textbf{1} & \textbf{2} & \textbf{3} & \textbf{4} & \textbf{5} & \textbf{6} & \textbf{7} & \textbf{8} & \textbf{9} & \textbf{10} & \textbf{11} & \textbf{12} \\
\hline
\textbf{A} & \textbf{A} & \textbf{A} & \textbf{A} & \textbf{A} & \textbf{A} & \textbf{A} & \textbf{A} & \textbf{A} & \textbf{A} & \textbf{A} & \textbf{A} \\
\hline
\end{tabular}
\end{center}

\vspace{1.5mm}
\noindent
\textbf{2. Phần 2: Trắc nghiệm Đúng -- Sai (2,0 điểm)}
\begin{itemize}
  \item \textbf{Câu 13 (1,0 điểm):} a) Đúng \quad b) Đúng \quad c) Sai \quad d) Đúng.
  \item \textbf{Câu 14 (1,0 điểm):} a) Đúng \quad b) Đúng \quad c) Sai \quad d) Đúng.
\end{itemize}
\textit{Quy tắc chấm: Đúng 1 ý được 0,1đ; đúng 2 ý được 0,25đ; đúng 3 ý được 0,5đ; đúng 4 ý được 1,0đ.}

\vspace{1.5mm}
\noindent
\textbf{3. Phần 3: Trắc nghiệm trả lời ngắn (2,0 điểm -- Mỗi câu đúng 0,5 điểm)}
\begin{itemize}
  \item \textbf{Câu 15:} $x \ge -4$.
  \item \textbf{Câu 16:} $4$.
  \item \textbf{Câu 17:} $60{,}6\text{ m}$ (hoặc $60{,}6$).
  \item \textbf{Câu 18:} $56{,}5\text{ m}^2$ (hoặc $56{,}52$).
\end{itemize}

\vspace{2mm}
\noindent
\textbf{\large B. LỜI GIẢI CHI TIẾT VÀ BAREM ĐIỂM TỰ LUẬN (3,0 ĐIỂM)}

\vspace{1.5mm}
\noindent
\textbf{Câu 19 (1,0 điểm).}
\begin{itemize}
  \item[a)] \textit{(0,5 điểm)} Hệ phương trình: $\begin{cases} 5x - 2y = 8 \quad (1) \\ 3x + 4y = 10 \quad (2) \end{cases}$.\\
  Nhân (1) với 2: $\begin{cases} 10x - 4y = 16 \\ 3x + 4y = 10 \end{cases}$ \dotfill \textbf{(0,25đ)}\\
  Cộng hai phương trình vế theo vế: $13x = 26 \iff x = 2$.\\
  Thay $x = 2$ vào (1): $5(2) - 2y = 8 \iff 2y = 2 \iff y = 1$.\\
  Vậy hệ phương trình có nghiệm duy nhất $(x; y) = (2; 1)$. \dotfill \textbf{(0,25đ)}

  \item[b)] \textit{(0,5 điểm)} Với $a > 0, a \ne 1$, ta có:
  \[
  \dfrac{\sqrt{a}+1}{\sqrt{a}-1} - \dfrac{\sqrt{a}-1}{\sqrt{a}+1} = \dfrac{(\sqrt{a}+1)^2 - (\sqrt{a}-1)^2}{a - 1} = \dfrac{(a + 2\sqrt{a} + 1) - (a - 2\sqrt{a} + 1)}{a - 1} = \dfrac{4\sqrt{a}}{a - 1} \quad \text{\textbf{(0,25đ)}}
  \]
  Do đó:
  \[
  M = \dfrac{4\sqrt{a}}{a - 1} : \dfrac{4\sqrt{a}}{a - 1} = \dfrac{4\sqrt{a}}{a - 1} \cdot \dfrac{a - 1}{4\sqrt{a}} = 1 \quad \text{\textbf{(0,25đ)}}
  \]
\end{itemize}

\vspace{1.5mm}
\noindent
\textbf{Câu 20 (1,5 điểm).}
\begin{itemize}
  \item[a)] \textit{(0,75 điểm)}
  Xét tam giác $BEC$ vuông tại $E$ và tam giác $BDC$ vuông tại $D$.\\
  Hai điểm $E, D$ cùng nhìn đoạn $BC$ dưới góc vuông nên $B, E, D, C$ cùng thuộc đường tròn đường kính $BC$ (đpcm). \dotfill \textbf{(0,25đ)}\\
  Vì $AK$ là đường kính của $(O)$ nên $\widehat{ACK} = 90^\circ \implies CK \perp AC$.\\
  Lại có $BH \perp AC$ (đường cao $BD$), suy ra $BH \parallel CK$. \dotfill \textbf{(0,25đ)}\\
  Tương tự, $\widehat{ABK} = 90^\circ \implies BK \perp AB$. Lại có $CH \perp AB$ (đường cao $CE$), suy ra $CH \parallel BK$.\\
  Tứ giác $BHCK$ có các cặp cạnh đối song song nên là hình bình hành (đpcm). \dotfill \textbf{(0,25đ)}

  \item[b)] \textit{(0,75 điểm)}
  Vì $BHCK$ là hình bình hành nên hai đường chéo $HK$ và $BC$ cắt nhau tại trung điểm của mỗi đường.\\
  Mà $M$ là trung điểm của $BC$, suy ra $M$ cũng là trung điểm của $HK$. Do đó ba điểm $H, M, K$ thẳng hàng (đpcm). \dotfill \textbf{(0,25đ)}\\
  Xét tam giác $AHK$: có $O$ là trung điểm của $AK$, $M$ là trung điểm của $HK$.\\
  Suy ra $OM$ là đường trung bình của tam giác $AHK$. \dotfill \textbf{(0,25đ)}\\
  Do đó $OM = \dfrac{1}{2} AH$, hay $AH = 2 OM$ (đpcm). \dotfill \textbf{(0,25đ)}
\end{itemize}

\vspace{1.5mm}
\noindent
\textbf{Câu 21 (0,5 điểm).}
\begin{itemize}
  \item[a)] \textit{(0,25 điểm)}
  Thể tích hình trụ là $V = \pi r^2 h = 54\pi \implies h = \dfrac{54}{r^2}$.\\
  Diện tích toàn phần lon nước ngọt:
  $S_{tp} = 2\pi r^2 + 2\pi r h = 2\pi \left( r^2 + \dfrac{54}{r} \right) = 2\pi \left( r^2 + \dfrac{27}{r} + \dfrac{27}{r} \right)$.\\
  Áp dụng BĐT Cô-si cho ba số dương:
  $r^2 + \dfrac{27}{r} + \dfrac{27}{r} \ge 3\sqrt[3]{r^2 \cdot \dfrac{27}{r} \cdot \dfrac{27}{r}} = 3 \cdot 9 = 27$.\\
  Suy ra $S_{tp} \ge 54\pi\text{ (cm}^2)$. Dấu đẳng thức xảy ra khi $r^2 = \dfrac{27}{r} \iff r^3 = 27 \iff r = 3\text{ cm}$.\\
  Khi đó $h = \dfrac{54}{3^2} = 6\text{ cm}$. Vậy diện tích nhỏ nhất khi $r = 3\text{ cm}$ và $h = 6\text{ cm}$. \dotfill \textbf{(0,25đ)}

  \item[b)] \textit{(0,25 điểm)}
  Áp dụng kỹ thuật Cauchy ngược dấu với $y > 0$: $\dfrac{x}{y^2 + 1} = x - \dfrac{xy^2}{y^2 + 1} \ge x - \dfrac{xy^2}{2y} = x - \dfrac{xy}{2}$.\\
  Tương tự: $\dfrac{y}{z^2 + 1} \ge y - \dfrac{yz}{2}$ và $\dfrac{z}{x^2 + 1} \ge z - \dfrac{zx}{2}$.\\
  Cộng vế theo vế: $S \ge (x + y + z) - \dfrac{xy + yz + zx}{2}$.\\
  Mà $xy + yz + zx \le \dfrac{(x+y+z)^2}{3} = \dfrac{3^2}{3} = 3 \implies S \ge 3 - \dfrac{3}{2} = \dfrac{3}{2}$.\\
  Dấu đẳng thức xảy ra khi $x = y = z = 1$. Giá trị nhỏ nhất của $S$ bằng $\dfrac{3}{2}$. \dotfill \textbf{(0,25đ)}
\end{itemize}

\end{document}
